\section{Conclusion and Future Work}
\label{sec:conclusion}

\paragraph{Answer to RQ1.} On the same hosted model, an agent harness that plans retrieval per bail factor, checks every claim and spars with a simulated prosecutor writes better-grounded bail briefs than the simpler systems we compared it with. The first full system led on objection recall and checked citations on development orders and trailed some baselines on the unsupported-fact rate; reading its outputs traced those losses to four causes, and four opt-in switches fixed them. On 80 sealed test orders the resulting system, c1457, beats the baselines on 46 comparisons and loses 5. It cites real authorities where the comparable retrieval baselines fabricate them, states fewer unsupported facts than any baseline, answers more of the prosecution's objections than the first full system, and needs 41\% fewer store calls than querying every store. One registered hypothesis, H-A1, is not supported: c1457 has one loss beyond the four declared in advance, the count of unsupported sentences against B5, which writes fewer sentences. No brief left the system without counsel's approval, and all 320 export attempts before Gate~2 were refused. The lesson we draw is methodological: the harness only helped once each of its decisions was measured against a simpler alternative and its losses were read case by case.

\paragraph{What was built.} Six agents with an event log; ten typed stores on four shelves; a page-level ingestion ladder for ten input kinds; evidence contracts for twelve bail factors and the orchestrator Sutradhar; claim checking and repair; prosecution sparring; truth maintenance with byte-identical redrafting; batch triage; the Rights Radar (all seven rules matched against the statute store by script); both approval gates and the pseudonymisation boundary; a web interface in English and Hindi and an API; six recorded demo cases with guided tours that go from the papers to a real Gate~2 approval and the exported application; and an opt-in sparring reserve that judges each reply right after it is drafted (on in the product at 200{,}000 tokens; Section~\ref{sec:limitations} gives its measured status). For evaluation: test sets T1 to T9, gold G1 to G7 without manual annotation, baselines B0 to B6, 46 opt-in switches for literature variants (61 of 61 arms pass an offline end-to-end audit), a scoring pipeline with paired bootstrap and Holm correction, and spend ledgers that cap every paid call. The test suite has about 5,000 tests.

\paragraph{Future work.} The next steps, in order of expected value, are a rerun of the sealed slice with the sparring reserve and the statute fixes, a reserve per case, quantity bands for NDPS from the charge sheet, several runs per system, and a second model family; Section~\ref{sec:limitations} gives the full list.

\paragraph{Code.} All code, tests, documents, the report source, run manifests and the tags \texttt{prereg-v1} and \texttt{eval-freeze-v1} are in \url{https://github.com/vishakkashyapk30/lma-group-project}. Data, model weights and keys are not in it; \texttt{make setup}, \texttt{make data} and \texttt{make models} rebuild them (Appendix, Reproducibility).
